% GENERATED FILE. Edit canonical lessons, then run npm run generate:books.
% canonical-source-hash: fnv1a64:452fa56331c6a5d7
% canonical-lessons: GU-C485-opresan, GU-C485-dakhal, GU-C485-tapkavvu, GU-C485-matra, GU-C485-pareji

\chapter{A Stay in Hospital}
\label{ch:gu-a-stay-in-hospital}
\hlchaptermodality{485}

\begin{chapteropening}
\textbf{By the end of this chapter:} \emph{I can talk about a hospital stay.}

Complete the last lesson of chapter 485: I can talk about a hospital stay.
\end{chapteropening}

\section[\texorpdfstring{\gu{ઓપરેશન} (opreśan)}{ઓપરેશન (opreśan)}]{\gu{ઓપરેશન} (opreśan) — an operation}

\label{lesson:GU-C485-opresan}



\begin{quote}
Before the new one: say the Gujarati for to cure, then the Gujarati for an injection.
\end{quote}

\subsection*{You'll want to know: \gu{ઓપરેશન}}
\textbf{\gu{ઓપરેશન}} — \emph{opreśan} — \textquotedblleft{}an operation\textquotedblright{}.

\begin{quote}
\gu{દાદાનું} \gu{ઓપરેશન} \gu{સારી} \gu{રીતે} \gu{થયું}.
\end{quote}

\begin{quote}
\emph{dādānũ opreśan sārī rīte thayũ.}
\end{quote}

\begin{quote}
Grandfather's operation went well.
\end{quote}

A neuter noun.

\begin{cousinweb}
\textbf{\gu{ઓપરેશન}} (\emph{opreśan}) is borrowed from English \emph{operation}.
\end{cousinweb}

\begin{grammarlens}[title={What you've built}]
One more word for describing what you lived and felt.
\end{grammarlens}

\subsection*{Your turn}
\begin{itemize}
\raggedright
  \item \emph{Say it:} \emph{opreśan}
  \item \emph{Say it:} \emph{opreśan}, once more
  \item \emph{Say it:} \emph{dādānũ opreśan sārī rīte thayũ.} — Grandfather's operation went well.
  \item \emph{Recall:} say the Gujarati for to cure, then the Gujarati for an injection, then say \emph{opreśan} again
\end{itemize}

\begin{tcolorbox}[breakable,colback=teal!4,colframe=teal!35!black,title={Before you move on}]
What does \gu{ઓપરેશન} mean? (\textquotedblleft{}An operation\textquotedblright{}.) Say it once more.
\end{tcolorbox}
\section[\texorpdfstring{\gu{દાખલ} (dākhal)}{દાખલ (dākhal)}]{\gu{દાખલ} (dākhal) — admitted (to hospital or school)}

\label{lesson:GU-C485-dakhal}



\begin{quote}
Before the new one: say the Gujarati for an injection, then the Gujarati for an operation.
\end{quote}

\subsection*{You'll want to know: \gu{દાખલ}}
\textbf{\gu{દાખલ}} — \emph{dākhal} — \textquotedblleft{}admitted (to hospital or school)\textquotedblright{}.

\begin{quote}
\gu{મીરા} \gu{બે} \gu{દિવસ} \gu{હોસ્પિટલમાં} \gu{દાખલ} \gu{હતી}.
\end{quote}

\begin{quote}
\emph{Mīrā be divas hospiṭalmā̃ dākhal hatī.}
\end{quote}

\begin{quote}
Mira was in hospital for two days.
\end{quote}

\begin{cousinweb}
\textbf{\gu{દાખલ}} (\emph{dākhal}) is usually traced to Persian \emph{dāxil}, \textquotedblleft{}entered, inside\textquotedblright{}, which comes from Arabic.
\end{cousinweb}

\begin{grammarlens}[title={What you've built}]
One more word for describing what you lived and felt.
\end{grammarlens}

\subsection*{Your turn}
\begin{itemize}
\raggedright
  \item \emph{Say it:} \emph{dākhal}
  \item \emph{Say it:} \emph{dākhal}, once more
  \item \emph{Say it:} \emph{Mīrā be divas hospiṭalmā̃ dākhal hatī.} — Mira was in hospital for two days.
  \item \emph{Recall:} say the Gujarati for an injection, then the Gujarati for an operation, then say \emph{dākhal} again
\end{itemize}

\begin{tcolorbox}[breakable,colback=teal!4,colframe=teal!35!black,title={Before you move on}]
What does \gu{દાખલ} mean? (\textquotedblleft{}Admitted (to hospital or school)\textquotedblright{}.) Say it once more.
\end{tcolorbox}
\section[\texorpdfstring{\gu{ટપકાવવું} (ṭapkāvvũ)}{ટપકાવવું (ṭapkāvvũ)}]{\gu{ટપકાવવું} (ṭapkāvvũ) — to put in drops}

\label{lesson:GU-C485-tapkavvu}



\begin{quote}
Before the new one: say the Gujarati for an operation, then the Gujarati for admitted.
\end{quote}

\subsection*{You'll want to know: \gu{ટપકાવવું}}
\textbf{\gu{ટપકાવવું}} — \emph{ṭapkāvvũ} — \textquotedblleft{}to put in drops\textquotedblright{}.

\begin{quote}
\gu{નર્સ} \gu{દિવસમાં} \gu{ત્રણ} \gu{વખત} \gu{આંખમાં} \gu{દવા} \gu{ટપકાવે} \gu{છે}.
\end{quote}

\begin{quote}
\emph{nars divasmā̃ traṇ vakhat ā̃khmā̃ davā ṭapkāve che.}
\end{quote}

\begin{quote}
Three times a day the nurse puts drops in my eye.
\end{quote}

Past, with \gu{તેણે} (he or she …): \gu{તેણે} \gu{આંખમાં} \gu{દવા} \gu{ટપકાવી}, \emph{teṇe ā̃khmā̃ davā ṭapkāvī}.

\begin{cousinweb}
\textbf{\gu{ટપકાવવું}} (\emph{ṭapkāvvũ}) is the causative of \textbf{\gu{ટપકવું}} (\emph{ṭapakvũ}), \textquotedblleft{}to drip\textquotedblright{}, built on a reconstructed \emph{*ṭapp-}, \textquotedblleft{}drip, drop\textquotedblright{}, an imitative word (CDIAL 5444).
\end{cousinweb}

\begin{grammarlens}[title={What you've built}]
One more word for describing what you lived and felt.
\end{grammarlens}

\subsection*{Your turn}
\begin{itemize}
\raggedright
  \item \emph{Say it:} \emph{ṭapkāvvũ}
  \item \emph{Say it:} \emph{ṭapkāvvũ}, once more
  \item \emph{Say it:} \emph{nars divasmā̃ traṇ vakhat ā̃khmā̃ davā ṭapkāve che.} — Three times a day the nurse puts drops in my eye.
  \item \emph{Recall:} say the Gujarati for an operation, then the Gujarati for admitted, then say \emph{ṭapkāvvũ} again
\end{itemize}

\begin{tcolorbox}[breakable,colback=teal!4,colframe=teal!35!black,title={Before you move on}]
What does \gu{ટપકાવવું} mean? (\textquotedblleft{}To put in drops\textquotedblright{}.) Say it once more.
\end{tcolorbox}
\section[\texorpdfstring{\gu{માત્રા} (mātrā)}{માત્રા (mātrā)}]{\gu{માત્રા} (mātrā) — a dose}

\label{lesson:GU-C485-matra}



\begin{quote}
Before the new one: say the Gujarati for admitted, then the Gujarati for to put in drops.
\end{quote}

\subsection*{You'll want to know: \gu{માત્રા}}
\textbf{\gu{માત્રા}} — \emph{mātrā} — \textquotedblleft{}a dose\textquotedblright{}.

\begin{quote}
\gu{દવાની} \gu{માત્રા} \gu{દાક્તરને} \gu{પૂછીને} \gu{લો}.
\end{quote}

\begin{quote}
\emph{davānī mātrā dāktarne pūchīne lo.}
\end{quote}

\begin{quote}
Ask the doctor about the dose before you take the medicine.
\end{quote}

A feminine noun.

\begin{cousinweb}
\textbf{\gu{માત્રા}} (\emph{mātrā}) is a tatsama, Sanskrit \emph{mātrā}, \textquotedblleft{}measure, quantity, size\textquotedblright{}, as Monier-Williams records it.
\end{cousinweb}

\begin{grammarlens}[title={What you've built}]
One more word for describing what you lived and felt.
\end{grammarlens}

\subsection*{Your turn}
\begin{itemize}
\raggedright
  \item \emph{Say it:} \emph{mātrā}
  \item \emph{Say it:} \emph{mātrā}, once more
  \item \emph{Say it:} \emph{davānī mātrā dāktarne pūchīne lo.} — Ask the doctor about the dose before you take the medicine.
  \item \emph{Recall:} say the Gujarati for admitted, then the Gujarati for to put in drops, then say \emph{mātrā} again
\end{itemize}

\begin{tcolorbox}[breakable,colback=teal!4,colframe=teal!35!black,title={Before you move on}]
What does \gu{માત્રા} mean? (\textquotedblleft{}A dose\textquotedblright{}.) Say it once more.
\end{tcolorbox}
\section[\texorpdfstring{\gu{પરેજી} (parejī)}{પરેજી (parejī)}]{\gu{પરેજી} (parejī) — a diet restriction}

\label{lesson:GU-C485-pareji}



\begin{quote}
Before the new one: say the Gujarati for to put in drops, then the Gujarati for a dose.
\end{quote}

\subsection*{You'll want to know: \gu{પરેજી}}
\textbf{\gu{પરેજી}} — \emph{parejī} — \textquotedblleft{}a diet restriction\textquotedblright{}.

\begin{quote}
\gu{ઓપરેશન} \gu{પછી} \gu{દાદાએ} \gu{તીખા} \gu{ભોજનની} \gu{પરેજી} \gu{પાળી}.
\end{quote}

\begin{quote}
\emph{opreśan pachī dādāe tīkhā bhojannī parejī pāḷī.}
\end{quote}

\begin{quote}
After the operation my grandfather kept off spicy food.
\end{quote}

A feminine noun.

\begin{cousinweb}
\textbf{\gu{પરેજી}} (\emph{parejī}) is usually traced to Persian \emph{parhēz}, \textquotedblleft{}abstinence, keeping away\textquotedblright{}.
\end{cousinweb}

\begin{grammarlens}[title={What you've built}]
One more word for describing what you lived and felt.
\end{grammarlens}

\subsection*{Your turn}
\begin{itemize}
\raggedright
  \item \emph{Say it:} \emph{parejī}
  \item \emph{Say it:} \emph{parejī}, once more
  \item \emph{Say it:} \emph{opreśan pachī dādāe tīkhā bhojannī parejī pāḷī.} — After the operation my grandfather kept off spicy food.
  \item \emph{Recall:} say the Gujarati for to put in drops, then the Gujarati for a dose, then say \emph{parejī} again
\end{itemize}

\begin{tcolorbox}[breakable,colback=teal!4,colframe=teal!35!black,title={Before you move on}]
What does \gu{પરેજી} mean? (\textquotedblleft{}A diet restriction\textquotedblright{}.) Say it once more.
\end{tcolorbox}
